% !TEX program = xelatex
\documentclass[10pt,a4paper]{article}

% ======================== 基础包 ========================
\usepackage[top=0.3cm,bottom=0.3cm,left=1.2cm,right=1.2cm]{geometry}
\usepackage{ctex}
\usepackage{titlesec}
\usepackage{enumitem}
\usepackage{xcolor}
\usepackage{hyperref}
\usepackage{fontawesome5}
\usepackage{tikz}

% ======================== 颜色定义 ========================
\definecolor{primary}{HTML}{2B4C7E}
\definecolor{darktext}{HTML}{2D2D2D}
\definecolor{graytext}{HTML}{666666}
\definecolor{rulecolor}{HTML}{B0B0B0}
\definecolor{tagbg}{HTML}{ECF0F5}
\definecolor{tagtext}{HTML}{2B4C7E}

% ======================== 超链接 ========================
\hypersetup{
    colorlinks=true,
    linkcolor=primary,
    urlcolor=primary,
    pdfauthor={徐其栋},
    pdftitle={徐其栋 - AI 应用 / 算法工程师}
}

% ======================== 字体设置 ========================
\setCJKmainfont{SimHei}[BoldFont=SimHei]
\setmainfont{Segoe UI}[BoldFont=Segoe UI Bold]

% ======================== 段落和间距 ========================
\setlength{\parindent}{0pt}
\setlength{\parskip}{0pt}
\pagestyle{empty}

% ======================== 标题样式 ========================
\titleformat{\section}
    {\large\bfseries\color{primary}}
    {}
    {0em}
    {}
    [\vspace{-0.5em}{\color{rulecolor}\rule{\textwidth}{0.6pt}}\vspace{-0.2em}]

\titlespacing*{\section}{0pt}{0.35em}{0.15em}

% ======================== 自定义命令 ========================
\newcommand{\projectblock}[3]{%
    {\bfseries\color{darktext}#1}\hfill{\color{graytext}#2}\\%
    {\small\color{graytext}#3}%
    \vspace{0.2em}%
}

% ======================== 文档开始 ========================
\begin{document}

% ======================== 头部 ========================
\begin{center}
    {\fontsize{26}{32}\selectfont\bfseries\color{primary} 徐其栋}

    \vspace{0.15em}

    {\color{graytext}
    \faPhone\hspace{0.3em}17268891140
    \hspace{1.0em}
    \faEnvelope\hspace{0.3em}\href{mailto:xqd922@outlook.com}{xqd922@outlook.com}
    \hspace{1.0em}
    \faGithub\hspace{0.3em}\href{https://github.com/xqd922}{github.com/xqd922}
    \hspace{1.0em}
    \faGlobe\hspace{0.3em}\href{https://1001.pp.ua}{个人主页}}

    \vspace{0.5em}

    {\color{darktext} 求职意向：\textbf{\color{primary}AI 应用工程师 / 算法工程师（初级） / AI 辅助开发}}
\end{center}

\vspace{0.2em}
{\color{rulecolor}\hrule height 0.8pt}

% ======================== 教育背景 ========================
\section{\faGraduationCap\hspace{0.4em}教育背景}

{\bfseries 滇西科技师范学院}\hfill{\bfseries 物联网工程\hspace{0.3em}|\hspace{0.3em}本科}

{\color{graytext} 2022.09 -- 2026.06}

\vspace{0.2em}

{\color{darktext}
\textbf{核心课程：}Python 程序设计、人工智能、图像处理、数据库应用、机器学习基础
}

% ======================== 专业技能 ========================
\section{\faCogs\hspace{0.4em}专业技能}

{\bfseries\color{primary}编程与数据：}\hspace{0.3em}{\color{darktext}Python、NumPy、Pandas、Matplotlib、scikit-learn}

\vspace{2pt}

{\bfseries\color{primary}机器学习与深度学习：}\hspace{0.3em}{\color{darktext}KNN、SVM、朴素贝叶斯、全连接神经网络、Keras / TensorFlow}

\vspace{2pt}

{\bfseries\color{primary}图像处理与 LLM：}\hspace{0.3em}{\color{darktext}OpenCV、OpenAI / Claude API、Prompt 工程、Claude Skills}

\vspace{2pt}

{\bfseries\color{primary}前端与工程：}\hspace{0.3em}{\color{darktext}Next.js、React、TypeScript、MDX、Vercel、Git、Linux、MySQL}

% ======================== 项目经历 ========================
\section{\faProjectDiagram\hspace{0.4em}项目经历}

% --- ML & DL 综合实战（主打项目） ---
\projectblock{Python 机器学习与深度学习综合实战}{2025.05 -- 2025.06}{技术栈：Python\hspace{0.3em}|\hspace{0.3em}scikit-learn\hspace{0.3em}|\hspace{0.3em}Keras\hspace{0.3em}|\hspace{0.3em}TensorFlow\hspace{0.3em}|\hspace{0.3em}NumPy / Pandas}

\begin{itemize}[leftmargin=1.5em, itemsep=2pt, parsep=0pt, topsep=2pt]
    \item 使用 scikit-learn 在 Iris 数据集上实现 KNN 与 SVM 双分类器，测试准确率均达 100\%。
    \item 独立从零实现朴素贝叶斯分类器并与 sklearn 库版本对比，熟悉精确率 / 召回率 / F1 多指标评估。
    \item 使用 Keras 搭建 128-64 双隐藏层全连接神经网络完成 MNIST 手写数字识别，覆盖数据预处理、训练、评估与可视化。
\end{itemize}

\vspace{0.15em}

% --- 个人博客 ---
\projectblock{AI 辅助开发的个人技术博客}{2024.11 -- 至今}{技术栈：Next.js 15\hspace{0.3em}|\hspace{0.3em}React 19\hspace{0.3em}|\hspace{0.3em}TypeScript\hspace{0.3em}|\hspace{0.3em}MDX\hspace{0.3em}|\hspace{0.3em}Vercel\hspace{0.3em}|\hspace{0.3em}Claude / GPT API}

\begin{itemize}[leftmargin=1.5em, itemsep=2pt, parsep=0pt, topsep=2pt]
    \item 独立使用 Next.js + TypeScript + MDX + Tailwind 开发个人博客，Vercel 持续部署，源码 \href{https://github.com/xqd922/blog}{github.com/xqd922/blog}，访问 \href{https://1001.pp.ua}{1001.pp.ua}。
    \item 以 Claude Code / ChatGPT 为结对开发助手，实际调用 OpenAI / Claude API 完成文本生成与代码辅助，熟悉 Prompt 工程与 Claude Skills 工作流。
\end{itemize}

\vspace{0.15em}

% --- 数字图像处理综合实践 ---
\projectblock{数字图像处理综合实践（车牌图像增强）}{2025.05}{技术栈：Python\hspace{0.3em}|\hspace{0.3em}OpenCV\hspace{0.3em}|\hspace{0.3em}NumPy\hspace{0.3em}|\hspace{0.3em}拉普拉斯锐化}

\begin{itemize}[leftmargin=1.5em, itemsep=2pt, parsep=0pt, topsep=2pt]
    \item 5 人小组完成车牌图像修复流水线，独立负责图像模糊改善模块，基于 OpenCV 实现拉普拉斯锐化核的二维卷积。
    \item 与中值滤波去噪、直方图均衡化提亮模块级联，形成完整图像增强流水线，熟悉 OpenCV 图像 I/O 与核卷积处理。
\end{itemize}

\vspace{0.15em}

% ======================== 个人评价 ========================
\section{\faUser\hspace{0.4em}个人评价}

2026 届本科生，求职方向 AI 应用 / 初级算法。熟悉 Python 数据科学栈，独立完成过经典机器学习、深度学习与图像处理项目；实际使用 OpenAI / Claude API 并熟练 Prompt 工程；具备 Next.js 前端开发能力，可快速进入 AI 应用研发工作。

\vspace{0.15em}

% ======================== 底部标签 ========================
\begin{center}
\tikz{\node[fill=tagbg, rounded corners=3pt, inner sep=5pt, text=tagtext, font=\small\bfseries] {Python};}
\hspace{0.2em}
\tikz{\node[fill=tagbg, rounded corners=3pt, inner sep=5pt, text=tagtext, font=\small\bfseries] {scikit-learn};}
\hspace{0.2em}
\tikz{\node[fill=tagbg, rounded corners=3pt, inner sep=5pt, text=tagtext, font=\small\bfseries] {Keras / TF};}
\hspace{0.2em}
\tikz{\node[fill=tagbg, rounded corners=3pt, inner sep=5pt, text=tagtext, font=\small\bfseries] {OpenCV};}
\hspace{0.2em}
\tikz{\node[fill=tagbg, rounded corners=3pt, inner sep=5pt, text=tagtext, font=\small\bfseries] {Claude / GPT API};}
\hspace{0.2em}
\tikz{\node[fill=tagbg, rounded corners=3pt, inner sep=5pt, text=tagtext, font=\small\bfseries] {Prompt 工程};}
\hspace{0.2em}
\tikz{\node[fill=tagbg, rounded corners=3pt, inner sep=5pt, text=tagtext, font=\small\bfseries] {Next.js / TS};}
\end{center}

\end{document}
